\scanpage{58}
\mspage{59}
\hfill n°269

voit que les formules
\[
(3)\qquad f(x)=f_1(x)-if_1(j(x))
\]
\[
(4)\qquad f(x)=f_2(j(x))+if_2(x)
\]
($x\in E$, $f\in E^*$, $f_1,f_2\in(E_0)^*$) établissent deux correspondances biunivoques entre $E^*$ et $(E_0)^*$.

3 - \uline{Formes bilinéaires associées à une forme hermitienne.}

Nous supposons encore ici que l'anneau $A$ est l'\uline{extension quadratique} $A=K(i)$ (où $i^2=-1$) d'un anneau commutatif $K$, et que $J$ est l'automorphisme $\lambda+i\mu\longrightarrow\lambda-i\mu$ de $A$ ($\lambda,\mu\in K$). Soient $E$ et $E'$ deux $A$-modules, $E_0$ et $E'_0$ les $K$-modules sous-jacents à $E$ et $E'$, $j$ et $j'$ les automorphismes $x\longrightarrow ix$ et $x'\longrightarrow ix'$ de $E$ et $E'$ (cf.n°2). Une forme $K$-bilinéaire $f$ sur $E_0\times E'_0$ sera dite \uline{invariante} par $j$ et $j'$ si l'on a
\[
(5)\qquad f(j(x),j'(x'))=f(x,x')
\]
pour $x\in E_0$ et $x'\in E'_0$. Remplaçant $x$ par $j(x)$, on voit que cette condition équivaut à
\[
(6)\qquad f(x,j'(x'))=-f(j(x),x')
\]
pour tous $x\in E_0$ et $x'\in E'_0$.

PROPOSITION 1 - \uline{Soit} $\Phi_1$ (resp. $\Phi_2$) \uline{une forme} $K$-\uline{bilinéaire sur} $E_0\times E'_0$ \uline{invariante par} $j$ \uline{et} $j'$. \uline{L'application qui à} $\Phi_1$ (resp. $\Phi_2$) \uline{fait correspondre l'application} $\Phi$ \uline{de} $E\times E'$ \uline{dans} $A$ \uline{définie par}
\[
(7)\qquad\Phi(x,x')=\Phi_1(x,x')+i\Phi_1(x,j(x'))
\]
\[
(\text{resp.}\quad(8)\qquad\Phi(x,x')=-\Phi_2(x,j(x'))+i\Phi_2(x,x')\ )
\]
($x,x'\in E$), \uline{est un isomorphisme du} $K$-\uline{espace vectoriel des formes} $K$-\uline{bilinéaires sur} $E_0\times E'_0$ \uline{invariantes par} $j$ \uline{et} $j'$ \uline{sur le}
\margin{57}
